\documentclass[12pt, titlepage,answers]{exam}
\usepackage{amsmath}
\usepackage{amsfonts}
\usepackage{amssymb}
\usepackage[english,bidi=basic]{babel}
\babelprovide[main,import,Alph=letters]{hebrew}
\babelfont[hebrew]{rm}{Hadasim CLM}
\babelfont[hebrew]{sf}{Miriam Mono CLM}
\usepackage{titlesec}
\title{מועד ב'}
\author{גלעד גור}
\begin{document}
\maketitle
\section*{הערה}
משום מה כתוב points הפוך, גם סדר הספרות הפוך אז אם כתוב $21$ הכוונה היא $12$ ולהיפך.\\
אם עשיתי טעויות בבקשה תכתבו לי - יונתן אבידור 054-2166594, אעדכן ואשלח מתוקן \\
מקווה שהלך לכולם מצויין ונתראה באלגברה ב'!
\begin{questions}
	\newpage
	\question
	תהא $T:\mathbb{R}_3\left[x\right] \to \mathbb{R}_3\left[x\right]$ המוגדרת כך
	\[
		T\left(ax^3+bx^2+cx+d\right) = \left(a+b+c\right)x^3 + \left(3a+c-d\right)x^2 + \left(a+b+c\right)x + \left(3a+c-d\right)
	\]
	\begin{parts}
		\part[12] הגדירו מהו פולינום אופייני עבור העתקה לינרארית וחשבו את הפולינום האופייני של $T$
		\begin{solution}
			מחשבים את המטריצה המייצגת לפי הבסיס הסטנדרטי $E=\left(x^3,x^2,x,1\right)$, יוצא
			\[
				\left[T\right]_E^E = \begin{pmatrix}
					1 & 1 & 1 & 0  \\
					3 & 0 & 1 & -1 \\
					1 & 1 & 1 & 0  \\
					3 & 0 & 1 & -1
				\end{pmatrix}
			\]
			סכום כל השורות זהה, לכן נעשה את הטריק שהיה גם במועד א'
			נעשה $C_1\to C_1 + C_2 + C_3 + C_4$ \\
			ואז כל האיברים ב-$C_1$ יהיו אותו דבר\\
			נעשה $R_2\to R_2 - R_1, R_3 \to R_3 - R_1, R_4 \to R_4 - R_1$
			ואז יהיה לנו $0$ בכל $C_1$ מלבד השורה הראשונה. כמה פעולות משם ויש לנו מטריצה משולשת עליונה ושלום על ישראל \\
			יוצאים ע"ע $3,0,0,-2$
		\end{solution}
		\part[5] האם $T$ הפיכה?
		\begin{solution}
			לא, $R_1 = R_3, R_2 = R_4$, לכן השורות ת"ל ו-$T$ לא הפיכה
		\end{solution}
		\part[13] האם $T$ לכסינה? אם כן מצאו בסיס $B$ ומטריצה אלכסונית $D$ כך ש $\left[T\right]^B_B =D$
		\begin{solution}
			טעיתי בזה אבל התשובה היא כן.\\
			בגלל ש-$r_a\left(\lambda\right) \geq r_g\left(\lambda\right) \geq 1$\\
			הריבוי האלגברי של כל הע"ע מלבד $0$ הוא $1$ ולכן בהכרח $r_g=r_a$\\
			אפשר לראות בעיניים שעבור ע"ע $0$ הדרגה היא $2$ ולכן מכיוון ש$4-2=2$ המימד העצמי של $0$ הוא $2$ וזה גם הריבוי האלגברי \\
			קיבלנו שהריבוי האלגברי זהה לריבוי הגיאומטרי לכל ע"ע וביחד זה יוצא $4=n$, לכסין.
			\[
				\begin{pmatrix}
					3 & 0 & 0 & 0  \\
					0 & 0 & 0 & 0  \\
					0 & 0 & 0 & 0  \\
					0 & 0 & 0 & -2
				\end{pmatrix}
			\]
			אין לי מושג מה הבסיס העצמי
		\end{solution}
	\end{parts}
	\newpage
	\question
	יהי $V$ מרחב וקטורי סוף-מימדי מעל שדה $\mathbb{C}$ ויהיו $T,S:V\to V$ העתקות לינאריות מתחלפות, כלומר $ST=TS$
	\begin{parts}
		\part[3] הראו כי ל-$T$ קיים לפחות ערך עצמי אחד $\alpha \in \mathbb{C}$
		\begin{solution}
			אם הפולינום האופייני פריק, יש ערך עצמי. כל פולינום מעל $\mathbb{C}$ הוא פריק.
		\end{solution}
		\part[7?] נסמן ב-$V_\alpha$ את המרחב העצמי עבור הערך העצמי $\alpha$ של $T$. הראו כי עבור כל $v\in V_\alpha$ מתקיים $S\left(v\right) \in V_\alpha$
		\begin{solution}
			\begin{gather*}
				T\left(v\right) = \alpha v \\
				\implies S\left(T\left(v\right)\right) = S\left(\alpha v\right) \\
				\implies T\left(S\left(v\right)\right) = \alpha S\left(v\right)
				\implies S\left(v\right) = V_\alpha
			\end{gather*}
		\end{solution}
		\part[3] תהא העתקה $S': V_\alpha  \to V_\alpha$ המוגדרת כ$S'\left(v\right)=S\left(v\right)$ עבור כל $v\in V_\alpha$. הראו כי ל-$S'$ קיים ערך עצמי.
		\begin{solution}
			אם הפולינום האופייני פריק, יש ערך עצמי. כל פולינום מעל $\mathbb{C}$ הוא פריק. \\
			שאלה טריקית, קצת לא קול לשים את אותו הדבר בתור שני סעיפים, למה שזה לא יהיה סעיף אחד בשווי 6 נקודות?
		\end{solution}
		\part[7?] הראו כי קיים וקטור עצמי משותף ל-$T$ ול-$S$
		\begin{solution}
			וואלה אין לי מושג, אני חירטטתי בטירוף
			כנראה קשור ל$S'\left(v\right)$ מהסעיף הקודם
		\end{solution}
	\end{parts}
	\newpage
	\question
	לפניכם שלוש טענות, קבעו אם הטענות נכונות או לא. אם נכונות יש להוכיח, אם לא יש להביא דוגמה נגדית מנומקת. בכל האחת מהשאלות $V$ הוא מרחב וקטורי סוף-מימדי מעל שדה $\mathbb{F}$ ו-$T$ היא העתקה לינארית $T:V\to V$.
	\begin{parts}
		\part[5] אם $Ker\left(T\right)=\mathrm{Im}\left(T\right)$ אזי $T$ היא העתקת האפס
		\begin{solution}
			לא, דוגמה נספק דוגמה נגדית
			\begin{gather*}
				\mathbb{F} = \mathbb{R}\\
				V=\mathbb{F}^2\\
				T\left(\begin{pmatrix}
						a \\b
					\end{pmatrix}\right)= \begin{pmatrix}
					0 \\a
				\end{pmatrix}
				Ker\left(T\right) = \begin{pmatrix}
					0 \\1
				\end{pmatrix}\\
				\mathrm{Im}\left(T\right) = \begin{pmatrix}
					0 \\1
				\end{pmatrix}\\
				T\left(\begin{pmatrix}
						1 \\1
					\end{pmatrix}\right) = \begin{pmatrix}
					0 \\1
				\end{pmatrix} \neq \begin{pmatrix}
					0 \\0
				\end{pmatrix}
			\end{gather*}
		\end{solution}
		\part[5] אם $T^2 = T$ אזי $Ker\left(T\right) \oplus \mathrm{Im}\left(T\right) = V$
		\begin{solution}
			נכון, צריך להראות שהם זרים ואז באמצעות שימוש בשני משפטי המימדים (שניהם רלוונטים כי $T$ אופרטור) הראות ש-$Ker+Im=V$
			על מנת להראות שהם זרים בוחרים איבר שאינו $0_V$ בגרעין, מניחים בשלילה שהוא בתמונה, לכן יש לו איבר מקור, מגיעים לסתירה שאומרת ש$v=0_V$ כסתירה להגדרת $v$
		\end{solution}
		\part[5] נניח $T$ הפיכה עם ע"ע $\lambda$. הוכיחו כי $\lambda \neq 0$ וכי $\lambda^{-1}$ ערך עצמי של $T^{-1}$ (אני חושב שהיה שם הרבה טקסט אז משהו לא מסתדר לי אבל זה מה שאני זוכר)
		\begin{solution}
			כן, אני טעיתי ממש בהוכחה, חירטטתי שטות לגבי הפולינום האופייני ומטריצה משולשת אבל תכלס
			$\lambda\neq 0$ כי אז יש שורת אפסים במטריצה האלכסונית שדומה ל-$A$ כש-$A$ היא מטריצה מייצגת כלשהי של $T$, דרגה נשמרת בדמיון ולכן גם ל-$A$ יש דרגה לא מלאה, סתירה לכך ש-$A$ הפיכה.\\
			על מנת להוכיח ש$\lambda^{-1}$ היא ע"ע של $A^{-1}$
			\begin{gather*}
				Av = \lambda v                                        \\
				\implies A^{-1}Av = A^{-1}\lambda v                   \\
				\implies v = \lambda A^{-1}v                          \\
				\implies \lambda^{-1}v = \lambda^{-1}\lambda A^{-1} v \\
				\implies \lambda^{-1}v = A^{-1}v
			\end{gather*}
			הלוואי והייתי חושב על זה במבחן פאק
		\end{solution}
	\end{parts}
	\newpage
	\question
	\begin{parts}
		\part[3] נסחו את משפט המימדים עבור העתקות לינאריות
		\begin{solution}
			יהי $\mathbb{F}$ שדה ו-$V,W$ מרחבים וקטורים ממימד סופי (יש שיגידו שרק $V$ צריך להיות ממימד סופי, אני כתבתי על שניהם כי אני לא יודע כלום על מרחבים ממימד אין-סופי) מעל $\mathbb{F}$ ותהא $T:V\to W$ העתקה לינארית, אזי
			\[
				\dim\left(Ker\left(T\right)\right) + \dim\left(\mathrm{Im}\left(T\right)\right) = \dim\left(V\right)
			\]
		\end{solution}
		\part[12?] הוכיחו את משפט המימדים עבור העתקות לינאריות
		\begin{solution}
			מסמנים את מימד התחום כ-$n$, מימד הגרעין כ-$k$ ורוצים להוכיח כי מימד התמונה הוא $n-k$\\
			בונים בסיס לגרעין, משלימים אותו לבסיס לתחום. התמונות של איברי קבוצה שפורשת את התחום תפרוש את התמונה ולכן נפעיל $T$ על הכל, נראה שהחלקים שמגיעים מהבסיס של הגרעין הם $0$ ולכן נבחר את ה$n-k$ וקטורים שנותרו ונטען כי הם בסיס לתמונה, על מנת לעשות את זה צריך להוכיח כי הם בתל \\
			עושים צ"ל שלהם, מלינאריות מוצאים שהם כולם בגרעין, משווים אותם לצ"ל של איברי בסיס הגרעין, מעבירים אגף ורואים שמדובר בצירוף של איברי בסיס התחום לכן מתאפס
		\end{solution}
	\end{parts}
	\newpage
	\question
	שאלת שיעורי הבית
	\begin{parts}
		\part[15] יהא $V$ ממימד $2$ מעל $\mathbb{R}$. תהא $T:V\to V$ העתקה לינארית המקיימת
		\[
			T^2 = -I
		\]
		הוכיחו כי קיים בסיס סדור $B$ ל-$V$ כך שמתקיים:
		\[
			\left[T\right]_B^B = \begin{pmatrix}
				0  & 1 \\
				-1 & 0
			\end{pmatrix}
		\]
	\end{parts}
	\begin{solution}
		יהא $0_V\neq v \in V$ \\
		שולפים מהמותן את הקבוצה $B=\left(T\left(v\right), v\right)$ היא קבוצה בגודל $2$ לכן היא בסיס אם היא בת"ל. \\
		מניחים בשלילה שהיא תלויה לינארית לכן
		\begin{gather*}
			\exists \alpha \in \mathbb{R}: v=\alpha T\left(v\right) \\
			\implies T\left(v\right) = \alpha T^2\left(v\right) \\
			\implies \alpha T\left(v\right) = -\alpha^2 v \\
			\implies v = -\alpha^2 v \\
			\implies \left(1+\alpha^2\right)v = 0\\
			\implies		1+\alpha^2 =0\\
			\alpha^2 = -1
		\end{gather*}
		סתירה, לא קיים סקלאר שכזה מעל $\mathbb{R}$. לכן הקבוצה הזו בת"ל ובסיס.\\
		עכשיו צריך להראות שהבסיס הזה אכן עונה לדרישה, מראים את המטריצה המייצגת ביחס לבסיס הזה
		\[
			\begin{aligned}
				T\left(T\left(v\right)\right) = -v = 0 \cdot T\left(v\right) -1\cdot v \\
				T\left(v\right) = 1 \cdot T\left(v\right) + 0 \cdot v
			\end{aligned} \implies \left[T\right]_B^B = \begin{pmatrix}
				0 & 1 \\ -1 & 0
			\end{pmatrix}
		\]
		איזה כיף
	\end{solution}
\end{questions}
\end{document}
